% Companion hardness section.
% Assumes amsmath, amssymb, amsthm, and the theorem/lemma/proof environments.
% No custom macros. Construction labels are prefixed hard:.
% The appended classical corollary uses the label full:np-complete.
% The complete-construction subsection can be moved to an appendix unchanged.

\section{Hardness parameterized by the deletion budget}
\label{hard:section}

The deletion budget and the preservation slack lead to different
parameterized complexity statements. The following companion result holds for the all-optional, unit-cost model. An original event is an
individually identified temporal contact; different contacts may have the
same endpoints and timestamp. Let \(E\) be the event set. For a retained
set \(S\subseteq E\), the quantity \(W_j(S)\) counts \(j\)-event walks
using events in \(S\), ending at the root, whose timestamps increase
strictly. Deleting \(D\) therefore leaves count \(W_j(E\setminus D)\).

\begin{theorem}[Hardness in the deletion budget]
\label{hard:budget}
Given an explicitly listed temporal in-arborescence of depth at most three,
deciding whether
\[
 |D|\le b,\qquad W_2(E\setminus D)\ge T_2,\qquad
 W_3(E\setminus D)\le T_3
\]
is \(\mathrm{W}[1]\)-hard parameterized by \(b\), even when every original
event is optional and has unit deletion cost. The restriction remains hard
with a single depth-one vertex and exactly one incoming depth-three event,
at timestamp zero, at each depth-two vertex that has such a predecessor.

The reduction from unweighted bipartite partial vertex cover with side
budgets \(a,q\) has \(b=a+2q\). Composed with the published reduction from
\(t\)-Clique, this gives \(b=t^2\).
\end{theorem}

\begin{proof}
The source is unweighted bipartite partial vertex cover with two separate
side budgets \cite[Definition~5 and Theorem~10]{mkrtchyan2024}.
The journal theorem states \(\mathrm{W}[1]\)-hardness for either side
budget; its displayed reduction from regular \(t\)-Clique uses
\(a=t\) and \(q=\binom{t}{2}\). Since both budgets are functions of
\(t\), that reduction also establishes hardness parameterized by
\(a+q\). The conference version is \cite[Theorem~6]{mkrtchyan2020}.
The regularization in the proof of Corollary~\ref{full:np-complete}
also preserves clique size, so the reduction may start from ordinary Clique.
The construction below has three steps.
Isolated rows and a pair of columns per source right vertex enforce the
two side budgets using the two temporal quotas.
Lemma~\ref{hard:retained-equivalence} proves this equivalence when the
middle events are temporarily required to survive.
Lemma~\ref{hard:replication} replaces that requirement by explicitly listed,
optional unit-cost events. Lemma~\ref{hard:size} proves polynomial encoded
size and the asserted parameter maps. These facts give a parameterized
many-one reduction and prove the theorem.
\end{proof}


The complete construction is in Appendix~\ref{hard:construction}.
Corollary~\ref{full:np-complete} gives NP-completeness of this decision
problem. It is a special case of the profile task: use unit costs, no
forced statuses, and select an entry with $w\le T_3$ and cost at most $b$.
The hardness parameter is $b$; no lower bound in preservation loss $R$ follows.
